% ═══════════════════════════════════════════════════════════════════════
%  Préambule propre au livre (PDF seulement)
% ═══════════════════════════════════════════════════════════════════════
% ── Coupure des URL ──────────────────────────────────────────────────
% Sans réglage, les longues adresses des notes de bas de page débordent
% dans la marge. xurl (chargé après hyperref, que la classe appelle en
% dernier) rend chaque caractère d'une URL sécable ; les trois lignes
% suivantes rétablissent une hiérarchie : couper de préférence après un
% « / », un « - », un « . » ou un « : », et à l'intérieur d'un mot
% seulement en dernier recours.
\usepackage{xurl}
\def\UrlBigBreaks{\do\:\do\/\do\-\do\.}
\mathchardef\UrlBigBreakPenalty=100
\mathchardef\UrlBreakPenalty=5000

% ═══════════════════════════════════════════════════════════════════════
%  Macros et styles TikZ du volume — partagés par le livre (PDF) et
%  les figures autonomes converties en SVG pour le web.
%  Toute macro utilisée dans une formule doit aussi être déclarée dans
%  macros-mathjax.tex (version web).
% ═══════════════════════════════════════════════════════════════════════
% ── Bibliothèques TikZ utilisées dans les figures ────────────────────
\usetikzlibrary{calc, arrows.meta, angles, quotes, patterns, positioning,
                decorations.pathreplacing, decorations.markings, fadings,
                backgrounds, fit}
% Angles, segments et points des figures en coordonnées polaires
\usepackage{tkz-euclide}
\tkzSetUpPoint[fill=black]
\definecolor{fill1}{RGB}{183,183,183}   % remplissage des angles

% ── Raccourcis du volume ─────────────────────────────────────────────
\providecommand{\piup}{\uppi}           % pi droit (anciennes figures)
\newcommand{\boitevide}{\mdlgwhtsquare} % case vide : 5 + □ = 3 (unicode-math)


% ── Chapitres d'exercices (séries 1 et 2) ─────────────────────────────
% Un chapitre entièrement consacré aux exercices (# Titre {.exercices}) :
% comme \sectionExercices, sans le titre « Exercices », puisque le
% chapitre en porte déjà un. Le filtre referme le fichier de solutions
% (\finExercices) avant le chapitre suivant.
\newcommand{\debutExercices}{%
  \Opensolutionfile{solutions}[sol-ch\thechapter]%
}
